
\medterm{U Wave} A small deflection following the T wave, best seen in precordial leads V2-V3. U waves
are thought to represent repolarization of the Purkinje fibers or papillary muscles.
Prominent U waves are associated with hypokalemia, bradycardia, and certain drugs
(digitalis, quinidine).

\medterm{UA} UA commonly means urinalysis, laboratory examination of urine for blood, protein, cells, glucose, infection, and other findings.

\medterm{Ubidecarenone} Coenzyme Q10, a lipid-soluble molecule involved in mitochondrial electron transport and cellular energy production; it is sold as a supplement and studied in selected cardiovascular and neuromuscular conditions.

\medterm{Ubiquinone} A lipid-soluble electron carrier in the mitochondrial membrane. See Coenzyme Q.

\medterm{Ubiquitin} A small regulatory protein that is covalently attached to other proteins to label them for degradation by the proteasome or to alter their cellular trafficking and signaling.

\medterm{Ubiquitination} Ubiquitination is attachment of ubiquitin to a target protein through an E1--E2--E3 enzyme cascade, marking it for proteasomal degradation or altering its localization, activity, or signaling.

\medterm{Ubiquitous} Present or found everywhere. In medical writing, the term is descriptive rather than a diagnosis and may refer to a substance, organism, or process that is widely distributed.

\medterm{UBT} UBT commonly means urea breath test, which detects active Helicobacter pylori infection by measuring labeled carbon dioxide after ingested urea.

\medterm{UDP-Glucuronosyltransferase} UDP-glucuronosyltransferase is an enzyme family that attaches glucuronic acid to bilirubin, medicines, hormones, and other compounds for elimination.

\medterm{Uhl Anomaly} A rare congenital absence or severe deficiency of right-ventricular myocardium, leaving
a thin-walled right ventricle. It may cause right-sided heart failure, arrhythmia,
cyanosis, or sudden deterioration in infancy or childhood.

\medterm{Ulcer} Full-thickness loss of epidermis extending into the dermis or deeper tissue, often
healing with a scar. Causes include vascular disease, pressure, infection, neuropathy,
inflammatory disease, trauma, and malignancy.

\medterm{Ulcer (Vascular)} A persistent break in the skin caused or maintained by inadequate blood flow, venous hypertension,
neuropathy, or related vascular disease.

\synonyms
Vascular

\medterm{Ulcer Healing Drugs} Medicines used to promote healing of gastrointestinal ulcers, usually by suppressing acid secretion or protecting the mucosa; treatment also addresses causes such as \textit{H. pylori} infection or NSAID use.

\medterm{Ulcer, Aphthous} Alternate index label; see \textit{Canker sore}.

\medterm{Ulcer, Peptic} Alternate index label; see \textit{Peptic ulcer}.

\medterm{Ulcer, Rectal} Alternate index label; see \textit{Solitary rectal ulcer syndrome}.

\medterm{Ulceration} The formation of an open sore caused by loss of surface tissue, usually involving skin or a mucous membrane and sometimes extending into deeper tissue.

\medterm{Ulcerative Colitis} A chronic inflammatory disease of the colon and rectum, usually causing continuous mucosal inflammation
that begins in the rectum and extends proximally. Symptoms may include bloody diarrhea,
urgency, abdominal pain, and extraintestinal manifestations.

\synonyms
Colitis, Ulcerative
Colitis Ulcerative (Ulcerative Colitis)

\medterm{Ulcerative Keratitis} An ulcerative infection or inflammatory defect of the cornea that can cause pain, redness, photophobia, and impaired vision and may threaten sight if untreated.

\medterm{Ulcerative Proctitis} Inflammation and ulceration limited to the lining of the rectum, often causing rectal bleeding, urgency, tenesmus, mucus, and diarrhea.

\medterm{Ullrich Congenital Muscular Dystrophy} An inherited muscle disease caused by problems with collagen VI, a protein that supports muscle and connective tissue. It begins in infancy with weakness, low muscle tone, tight joints near the body, unusually loose joints in the hands and feet, and later breathing problems.

\synonyms
UCMD; Ullrich Scleroatonic Muscular Dystrophy

\medterm{Ulna} Medial forearm bone, longer than the radius. Proximal end: trochlear notch, olecranon, coronoid process. Shaft. Distal end: head, styloid process. The trochlear notch articulates with the trochlea forming the primary elbow hinge.

\medterm{Ulnar} Relating to the ulna bone of the forearm or the ulnar side (medial side) of the hand.
Ulnar nerve compression causes paresthesia in medial one and a half fingers.

\medterm{Ulnar Artery} The larger terminal branch of the brachial artery at the cubital fossa, running medially
down the forearm to the wrist (palpable medial to the pisiform) and hand. Supplies the
medial forearm and ulnar side of the hand; forms the superficial palmar arch. Provides
the common interosseous artery.

\medterm{Ulnar Vein} Deep veins accompanying the ulnar artery that drain the medial forearm and join the radial veins to form the brachial veins.

\medterm{Ulnar Hypoplasia} Underdevelopment or partial absence of the ulna, usually presenting as a congenital
forearm and hand anomaly. It may cause forearm shortening, radial deviation of the hand,
limited elbow or wrist movement, and associated digital or systemic abnormalities.

\medterm{Ulnar Nerve} From medial cord (C8-T1), passing behind the medial epicondyle (cubital tunnel), into
the forearm innervating flexor carpi ulnaris and medial FDP half. In the hand innervates
hypothenar muscles, interossei, medial two lumbricals, and adductor pollicis. Sensation
to medial 1.5 digits.

\medterm{Ulnar Nerve Dysfunction} Impairment of the ulnar nerve, commonly caused by compression at the elbow (cubital
tunnel) or wrist (Guyon canal). Symptoms include numbness or tingling in the little and
ring fingers, weakened grip, and hand-muscle wasting; severe lesions may produce an
ulnar claw.

\medterm{Ulnar Nerve Entrapment} Compression or irritation of the ulnar nerve, commonly at the elbow or wrist, causing numbness or tingling in the ring and little fingers and sometimes weakness of the hand.

\medterm{Ulnar Paradox} The clinical observation that a lesion of the ulnar nerve at the wrist may cause more clawing of the fingers than a lesion higher in the forearm, because the ulnar-innervated finger flexors remain intact.

\medterm{Ulnar Wrist Pain} Pain on the side of the wrist nearest the little finger, often involving the ulna or nearby tendons and ligaments.

\synonyms
Ulnar-Sided Wrist Pain

\medterm{Ulnar-Sided Wrist Pain} Alternate index label; see \textit{Ulnar wrist pain}.

\medterm{Ulocladium} A genus of dematiaceous filamentous fungi found in soil and plant material. It is an uncommon opportunistic pathogen that can cause keratitis, sinusitis, or invasive infection in susceptible people.

\medterm{Ultracentrifugation} Ultracentrifugation separates particles or macromolecules at very high centrifugal force according to size, shape, and density; it is used to study viruses, organelles, lipoproteins, and protein complexes.

\medterm{Ultradian Rhythm} A biological rhythm that repeats more than once in 24 hours, such as pulsatile hormone secretion, sleep-stage cycling, or periodic changes in heart and breathing patterns.

\medterm{Ultrafiltration} Separation in which pressure pushes water and small dissolved substances through a membrane while larger particles stay behind. The process helps form urine and is also used during dialysis.

\medterm{Ultralente Insulin} Ultralente insulin is a prolonged-acting zinc insulin suspension formerly used for basal insulin replacement; it has largely been replaced by newer long-acting analogues and can cause hypoglycaemia and dosing variability.

\medterm{Ultrasonic Therapy} Use of high-frequency sound energy for therapeutic purposes such as selected physical-therapy applications or tissue ablation; benefits and risks depend on the device and treatment site.

\medterm{Ultrasonography} A scan using high-frequency sound waves to create images of organs, tissues, vessels, or a fetus without radiation.

\medterm{Ultrasound} An imaging test that uses high-frequency sound waves to show organs, blood vessels, tissues, or a developing fetus without ionizing radiation. Different examinations can assess the abdomen, pelvis, kidneys, thyroid, heart, blood flow, or testicles.

\medterm{Ultrasound Biomicroscopy} High-frequency ultrasound imaging used to visualize the anterior segment of the eye, including the cornea, iris, ciliary body, lens, and angle when ordinary examination or imaging is insufficient.

\medterm{Ultrasound Contrast} Tiny gas-filled bubbles given into a vein to improve ultrasound images of blood flow and selected organs. They are usually cleared through the lungs but can rarely cause allergic or cardiopulmonary reactions.
\medterm{Ultrasound Therapy} Use of sound waves at therapeutic intensity to heat or mechanically affect tissue. It is used in selected rehabilitation settings, but evidence and suitability vary by condition and it is not appropriate over every body area.

\medterm{Ultrasound Transducer} A device that emits and receives high-frequency sound waves to generate diagnostic
images or guide procedures. Linear, curvilinear, phased-array, and endocavitary
transducers are selected according to depth and anatomic application.

\medterm{Ultraviolet Microscopy} Microscopy that uses ultraviolet wavelengths to produce images or fluorescence from structures and substances that absorb or emit ultraviolet light. It has specialized laboratory and research applications.

\medterm{Ultraviolet Rays} Invisible electromagnetic radiation with wavelengths shorter than visible violet light. Excess exposure can cause sunburn, photoaging, eye injury, and skin cancers; controlled ultraviolet exposure has selected therapeutic uses.

\medterm{Ultraviolet Therapy} Controlled exposure to ultraviolet light used for selected skin disorders such as psoriasis, eczema, and vitiligo; treatment requires eye and skin protection and monitoring for burns and long-term skin cancer risk.

\medterm{Umbilical} Relating to the umbilicus or navel, including structures, lesions, discharge, or abnormalities located at or connected with it.

\medterm{Umbilical Artery Doppler} Ultrasound assessment of blood flow in the umbilical arteries to evaluate placental
resistance. Normal flow shows a pulsatility index decreasing throughout pregnancy with
forward diastolic flow. Absent or reversed end-diastolic flow indicates severe placental
dysfunction and is associated with increased perinatal morbidity and mortality.

\medterm{Umbilical Catheters} Umbilical catheters provide vascular access through a newborn’s umbilical vessels for fluids, medicines, monitoring, or blood sampling while complications are monitored.

\medterm{Umbilical Cord} The connection between the fetus and the placenta, containing two umbilical arteries
(carrying deoxygenated blood from the fetus to the placenta) and one umbilical vein
(carrying oxygenated blood from the placenta to the fetus). The vessels are embedded in
Wharton jelly, a mucopolysaccharide matrix that prevents compression.

\medterm{Umbilical Cord Blood Gas} Analysis of arterial and venous cord blood collected after delivery. The umbilical artery
best reflects fetal acid-base status, while the umbilical vein reflects placental blood. Results
are interpreted using gestational age, delivery circumstances, pH, carbon dioxide, and base
deficit or excess.

\medterm{Umbilical Cord Prolapse} Descent of the umbilical cord through the cervix alongside or past the presenting fetal
part, occurring when the membranes rupture. Risk factors: malpresentation (transverse
lie, breech), preterm/low birth weight, polyhydramnios, multiparity, and artificial ROM
with high presenting part.

\medterm{Umbilical Granuloma} A small moist pink or red lump of healing tissue at a baby's navel after the cord falls off. It may ooze clear fluid but should not cause spreading redness, fever, or marked tenderness; a clinician can treat it and exclude infection or another abnormal connection.

\synonyms
Granuloma Umbilicale; Neonatal Umbilical Granuloma

\medterm{Umbilical Hernia} A hernia through the umbilical ring, common in infants (most umbilical hernias close
spontaneously by age 4-5) and adults (risk factors: obesity, pregnancy, ascites,
cirrhosis). Presents as a reducible bulge at the umbilicus, worse with increased
intra-abdominal pressure. In adults: risk of incarceration and strangulation is higher.

\medterm{Umbilical Region} Central region around the umbilicus containing small intestine, transverse colon, and
great vessels. The umbilicus is a landmark for assessing ascites, caput medusae in
portal hypertension, and umbilical hernias. Percussion and auscultation here are
essential components of the abdominal examination.

\medterm{Umbilical Vein} The umbilical vein carries oxygen- and nutrient-rich blood from the placenta to the fetus during pregnancy.

\medterm{Umbilication} Formation of a central hollow or navel-like depression in a skin lesion or other body structure. It can occur in some infections, inflammatory lesions, tumors, or normal anatomical features and must be interpreted with the lesion’s appearance.
\medterm{Umbilicus} The navel, a scar on the anterior abdominal wall marking the former attachment of the umbilical cord. It can be affected by hernia, infection, discharge, or metastatic disease.

\medterm{Umbilicus (Navel)} The scar on the anterior abdominal wall marking the site of attachment of the umbilical
cord in fetal life, located at the level of the L3-L4 intervertebral disc, approximately
midway between the xiphoid process and the pubic symphysis.

\synonyms
Navel

\medterm{Uncertainty} Lack of complete certainty about a diagnosis, prognosis, measurement, or treatment effect. Clinical uncertainty is managed by combining available evidence with follow-up, additional testing, and shared decision-making.

\medterm{Uncinate Fit} A nonstandard or unclear historical label in the source dictionary; its intended clinical meaning should be verified against the original reference.

\medterm{Uncompetitive Inhibition} A type of enzyme inhibition where the inhibitor binds only to the enzyme-substrate
complex at a site distinct from the active site, forming an inactive ESI complex. The
inhibitor cannot bind to the free enzyme.

\medterm{Unconscious} Not awake and not responsive to ordinary external stimuli. Unconsciousness may result from syncope, seizure, trauma, intoxication, metabolic disease, stroke, or other emergencies.

\medterm{Unconsciousness} Loss of awareness and normal response to surroundings, ranging from brief fainting to prolonged coma. It can result from reduced brain blood flow, injury, seizures, low blood sugar, drugs, or serious illness.

\medterm{Unconsciousness, Temporary} A brief episode of loss of awareness and responsiveness, often caused by syncope but also possible with seizure, concussion, hypoglycemia, or other conditions requiring assessment.

\medterm{Underventilation} Breathing that is too shallow or slow to remove enough carbon dioxide, causing hypercapnia and sometimes low oxygen; causes include sedatives, neuromuscular weakness, obesity hypoventilation, and lung disease.

\medterm{Underweight} A body weight below the usual range for height. In adults, a Body Mass Index below 18.5 kg/m\(^2\) is commonly classified as underweight; illness or poor nutrition may contribute.

\medterm{Undescended Testicle} A testicle that has not moved into the scrotum before birth, also called cryptorchidism. It may descend during early infancy, but a testicle that remains high can have increased risks of infertility, twisting, injury, and cancer and should be assessed.

\synonyms
Testicle, Undescended

\medterm{Undescended Testicle Repair} Surgical relocation and fixation of a testis that has not descended into the scrotum, usually performed in childhood to preserve function and reduce complications.

\medterm{Undescended Testis} Failure of one or both testes to descend into the scrotum. Some descend during early
infancy, but persistent cryptorchidism can impair fertility and increase the risk of torsion, hernia, or testicular
cancer.

\medterm{Undescended Testis (Cryptorchidism)} Congenital failure of one or both testes to descend along the normal path
into the scrotum. Persistent cases require evaluation because they may affect fertility and increase the risk of torsion,
hernia, or testicular cancer.

\synonyms
Cryptorchidism

\medterm{Undetermined} Not established or not able to be classified with the available clinical information. An undetermined result may require repeat assessment or additional testing.

\medterm{Undifferentiated} Lacking recognizable specialized features. In pathology, an undifferentiated tumor has cells that cannot be assigned readily to a specific tissue of origin based on their appearance.

\medterm{Undifferentiated Cancer} A cancer whose cells lack recognizable specialized features, making the tissue of origin or exact subtype difficult to identify; diagnosis requires pathology and staging.

\medterm{Undifferentiated Carcinoma} A high-grade carcinoma whose cells lack recognizable features of a specific tissue or tumor subtype; diagnosis and treatment require pathology, immunostaining, and staging.

\medterm{Undifferentiated Pleomorphic Sarcoma} An uncommon aggressive cancer of soft tissue or bone whose cells have no clearly identifiable line of differentiation.

\synonyms
Malignant Fibrous Histiocytoma

\medterm{Undifferentiated Pleomorphic Sarcoma Of The Breast} A rare high-grade breast sarcoma formerly called malignant fibrous histiocytoma. It is a diagnosis of exclusion after other specific sarcomas and metaplastic carcinoma have been ruled out histologically.
\medterm{Undulant Fever} A zoonotic infectious disease caused by gram-negative coccobacilli of the genus
Brucella, transmitted to humans through contact with infected animals (cattle, goats,
pigs, sheep) or consumption of unpasteurized dairy products. See Brucellosis.

\medterm{Undulate} To move or have a wave-like contour. In a clinical description, it may refer to a palpable fluid wave, a wavy tissue margin, or a rhythmic motion.

\medterm{Undulate Ribs} Ribs with a wavy or irregular contour, a descriptive anatomic finding that may reflect a developmental variation, deformity, prior injury, or underlying bone disease.


\medterm{Unequivocal} Clear and not open to more than one reasonable interpretation. An unequivocal clinical finding or test result provides strong evidence for the stated interpretation when considered in context.

\medterm{Unerupted Tooth} A tooth that has not emerged into the oral cavity at the expected time. It may be normally developing, delayed, impacted, or prevented from erupting by obstruction or abnormal position.

\medterm{Unevaluable} Unable to be assessed or interpreted reliably because of inadequate data, poor test quality, an unavailable endpoint, or another limiting circumstance.

\medterm{Ungual} Relating to a fingernail or toenail. The term may describe a finding beneath, around, or within the nail, such as a bruise, fungal infection, wart, or pigmented lesion.

\medterm{Unguentum} A semisolid, greasy preparation intended for application to the skin or, when formulated for that purpose, to mucous membranes or the eyes.

\medterm{Unicornuate Uterus} A congenital condition in which only one side of the uterus develops fully, usually with one fallopian tube. Some people have no symptoms, but it can increase risks of infertility, pregnancy outside the uterus, miscarriage, or premature birth and may require specialist pregnancy care.

\medterm{Unilaminar Epithelium} An epithelium consisting of a single layer of cells. Its structure is adapted for functions such as diffusion, absorption, secretion, or filtration depending on the organ.

\medterm{Unilateral} Affecting, occurring on, or involving one side of the body or one paired organ, such as unilateral weakness or unilateral eye disease.

\medterm{Unilateral Gynecomastia} Enlargement of breast glandular tissue affecting one breast in a male patient; a new, firm, fixed, painful, or nipple-associated mass requires clinical assessment.

\medterm{Unilateral Microphthalmos} A congenital condition in which one eye is abnormally small. It may be associated with other ocular malformations and can cause substantial visual impairment.

\medterm{Uninodular Goiter} A goiter containing one dominant thyroid nodule. Evaluation generally includes thyroid-function testing and imaging, with further assessment to determine whether the nodule is benign or malignant.

\medterm{Uniparental Disomy} Inheritance of both copies of a chromosome or chromosome segment from one parent and
none from the other. It can cause disease through abnormal genomic imprinting,
homozygosity for a recessive variant, or chromosomal dosage effects.

\medterm{Unipolar Neuron} A neuron with a single process extending from its cell body. True unipolar neurons occur in some invertebrates; many sensory neurons in humans are more accurately described as pseudounipolar.

\medterm{Unit} A standard amount used to express measurement, dose, concentration, or a healthcare service in an agreed system.


\medterm{Unopposed Estrogen} Estrogen treatment given without a progestogen to someone with a uterus. Long-term use can over-stimulate the uterine lining and increase cancer risk, so protection is usually recommended when appropriate.

\medterm{Unresectable} Describing a tumor or lesion that cannot be completely removed safely by surgery because of its location, extent, invasion of vital structures, or the patient’s operative risk.

\medterm{Unresolved Losses} Important losses whose effects have not been fully recognized or worked through. They may involve death, health, relationships, work, home, or identity and can contribute to lasting grief, anger, guilt, anxiety, or difficulty moving forward.

\medterm{Unresponsive Coma} A state of profound unconsciousness in which a person cannot be awakened or follow commands. Causes include brain injury, stroke, seizures, poisoning, infection, and metabolic failure; emergency evaluation is required.
\medterm{Unsaturated Phytosterol} A plant sterol containing one or more carbon--carbon double bonds. Phytosterols can reduce intestinal cholesterol absorption, although clinical effects depend on intake, formulation, and patient context.

\medterm{Unstable Angina} Myocardial ischemia at rest, new-onset severe angina, or a crescendo pattern (more
frequent, longer, lower-threshold) without evidence of myocardial necrosis (normal
troponin). Pathophysiology: acute plaque rupture with thrombosis causing transient
severe coronary obstruction.

\medterm{Unsteadiness} Impaired balance or gait stability, which may result from vestibular disease, neurologic disorders, weakness, medication effects, visual impairment, or musculoskeletal problems.

\medterm{Untidy Wound} An irregular or contaminated wound with crushing, devitalized tissue, tissue loss, foreign material, or associated injury to vessels, nerves, tendons, or bone. It has increased infection risk and may require exploration, irrigation, debridement, and delayed or staged closure.

\medterm{Untranslated Region} A part of an RNA molecule that is not translated into protein; the 5-prime and 3-prime untranslated regions regulate RNA stability, localization, and translation.

\medterm{Unwell} A general description of feeling ill or having reduced health, without identifying a specific disease. Persistent or severe symptoms require assessment for an underlying cause.

\medterm{Upamostat} An investigational small-molecule serine-protease inhibitor studied for inflammatory and infectious diseases. It is not an established routine treatment, and its clinical role depends on trial evidence and regulatory approval.

\medterm{Upper Airway Biopsy} Removal of tissue from the nasal cavity, pharynx, larynx, or related upper airway for microscopic evaluation of suspected infection, inflammation, infiltrative disease, or tumor.

\medterm{Upper Airway Obstruction} Narrowing of the airway above the carina, causing inspiratory stridor, dyspnea, and
characteristic changes on flow-volume loops (flattening of the inspiratory limb). Causes
include tumors, stenosis, goiter, and vocal cord dysfunction. It is a medical emergency
when severe.

\medterm{Upper Airway Resistance Syndrome} Increased effort to breathe through a narrowed upper airway, often causing sleep disruption and daytime sleepiness.

\synonyms
UARS

\medterm{Upper Arm} The portion of the upper limb between the shoulder and elbow, containing the humerus and muscles, nerves, and vessels that serve the arm.

\medterm{Upper Endoscopy} Upper endoscopy: flexible endoscopic examination of the esophagus, stomach, and duodenum. See Esophagogastroduodenoscopy (EGD).
\medterm{Upper Esophageal Sphincter} A ring of muscle at the top of the esophagus that opens during swallowing and coordinates food passage.

\synonyms
UES

\medterm{Upper Extremity} The anatomic region extending from the shoulder through the arm, forearm, wrist, and hand. Disorders may involve its bones, joints, muscles, nerves, vessels, or skin.

\medterm{Upper Extremity Venous Doppler} An ultrasound examination of the arm and related veins that assesses compressibility and blood flow, commonly to evaluate suspected venous thrombosis.

\medterm{Upper Gastrointestinal Bleeding} Bleeding from esophagus, stomach, or duodenum. Causes: peptic ulcer, varices,
Mallory-Weiss tear, gastritis. Presentation: hematemesis, melena, hematochezia.
Management: ABCs, IV access, type and crossmatch, PPI infusion, endoscopy within 24
hours. Variceal bleeding: octreotide, antibiotics, band ligation.

\medterm{Upper GI And Small Bowel Series} A fluoroscopic x-ray examination in which swallowed contrast outlines the esophagus, stomach, duodenum, and small intestine to evaluate structure and transit.

\medterm{Upper GI Series} Fluoroscopic examination of the esophagus, stomach, and duodenum after oral barium
contrast. Single-contrast: looks for gross abnormalities (obstruction, masses).
Double-contrast (barium + gas): better mucosal detail for ulcers, erosions, and early
malignancies.

\medterm{Upper Leg} The portion of the lower limb between the hip and knee, containing the femur and the thigh muscles, nerves, and blood vessels.

\medterm{Upper Limb Disorders} Conditions affecting the shoulder, arm, elbow, forearm, wrist, or hand, including musculoskeletal, nerve, vascular, and work-related disorders.

\medterm{Upper Motor Neuron} A neuron in the brain or spinal cord that influences lower motor neurons controlling skeletal muscle. Damage can cause weakness, spasticity, increased reflexes, and pathologic reflexes.

\medterm{Upper Motor Neuron Syndrome} A pattern caused by injury to corticospinal pathways above the anterior horn cell or
cranial motor nucleus. Findings may include weakness with spasticity, hyperreflexia,
clonus, and an extensor plantar response, although acute lesions can initially produce
flaccidity.

\medterm{Upper Respiratory Infection} An infection of the nose, sinuses, pharynx, or larynx, most often caused by a virus and producing symptoms such as congestion, sore throat, cough, or runny nose.

\medterm{Upregulation} An increase in the activity, amount, or expression of a gene, receptor, enzyme, or signaling pathway in response to developmental, physiologic, or disease-related signals.

\medterm{Upright Labor Positions} Maternal positions during labor, such as walking, standing, squatting, kneeling, or sitting upright.
They may improve comfort and mobility and can be used when safe and acceptable to the
laboring person and clinical team.

\medterm{Uptake} The accumulation of a substance, such as a radiotracer, inside cells or tissues. It depends on blood flow, transport across cell membranes, metabolism, binding, and the tissue’s normal or abnormal activity.

\medterm{Urachal Carcinoma} A rare cancer arising from a leftover connection between the bladder and the navel, usually near the top of the bladder. It may cause blood or mucus in the urine, pelvic pain, or urinary changes.

\medterm{Urachal Cyst} A cystic remnant of the embryonic urachus between the bladder and umbilicus. It may be
asymptomatic or become infected, causing lower-abdominal pain, fever, umbilical
discharge, or urinary symptoms.

\medterm{Urachal Mucinous Cystic Tumor} A rare tumor arising from the urachal remnant between the bladder dome and umbilicus. It
may be benign or malignant, with mucinous histology, and requires surgical excision.

\medterm{Urachus} An embryologic connection between the fetal bladder and umbilicus that normally becomes the median umbilical ligament after birth. Persistent urachal remnants can form cysts, sinuses, or other complications.

\medterm{Uracil} A pyrimidine nitrogenous base found principally in RNA. It pairs with adenine and is also a component of several metabolic intermediates and nucleotide medicines.

\medterm{Uracil Mustard} An older nitrogen-mustard alkylating compound related to uracil and studied or used historically as an antineoplastic drug. Its clinical use is limited by toxicity and the availability of newer treatments.

\medterm{Uracil Nucleotides} Nucleotides containing uracil, including UMP, UDP, and UTP. They participate in RNA synthesis, carbohydrate metabolism, and activation of sugars for glycoconjugate production.

\medterm{Uraemia} The clinical syndrome resulting from accumulation of urea and other nitrogenous waste
products in the blood in advanced chronic kidney disease (CKD stage 4-5) and acute
kidney injury.

\medterm{Uraemic Syndrome} The collection of symptoms and biochemical abnormalities caused by severe accumulation of uraemic waste products in advanced kidney failure.

\medterm{Uranium} A naturally occurring radioactive heavy metal. Internal exposure can damage the kidneys, while inhaled or retained radioactive material can also expose tissues to ionizing radiation.

\medterm{Urate} The ionized form of uric acid, the end product of purine metabolism in humans. Excess serum urate can promote monosodium urate crystal deposition and gout.

\medterm{Urate-Lowering Therapy} Long-term treatment that reduces serum urate to dissolve monosodium urate deposits and
prevent gout flares and structural damage. Xanthine oxidase inhibitors and uricosuric
medicines are selected according to renal function, comorbidity, and treatment goals.

\medterm{Urates} Salts or ions of uric acid. Monosodium urate crystals can accumulate in joints in gout, and urate crystals may contribute to urinary stones.

\medterm{Urea} The principal nitrogen-containing waste product of protein metabolism, formed in the liver and excreted mainly by the kidneys.

\medterm{Urea Breath Test} A noninvasive test for active Helicobacter pylori infection: the person drinks labeled urea, and labeled carbon dioxide in exhaled breath indicates bacterial urease activity.

\medterm{Urea Cycle} A metabolic pathway occurring in both the mitochondria and cytoplasm of hepatocytes that
converts toxic ammonia (NH3), derived from amino acid catabolism, into nontoxic,
water-soluble urea for excretion by the kidneys.

\medterm{Urea Cycle Disorder} A group of inherited disorders caused by deficiency of an enzyme or transporter in the urea cycle,
leading to impaired ammonia clearance. Ornithine transcarbamylase deficiency is an
X-linked form; presentation and severity vary by the gene and residual enzyme activity.

\medterm{Urea Nitrogen Urine Test} A urine test measuring urea nitrogen excretion, interpreted with blood urea nitrogen, creatinine, diet, and urine volume to assess protein metabolism and renal handling.

\medterm{Ureaplasma} A genus of very small bacteria without a cell wall that can colonize the genital or urinary tract and is sometimes associated with urethritis or pregnancy complications.

\medterm{Ureaplasma Urealyticum} A cell-wall-free bacterium that can colonize the genital or urinary tract and is sometimes associated with nongonococcal urethritis, pregnancy complications, or infection in vulnerable newborns.

\medterm{Uremia} A clinical syndrome caused by advanced kidney failure and accumulation of waste products and other biologically active
substances. Symptoms may include nausea, itching, fatigue, confusion, bleeding, or
pericarditis; laboratory values and symptoms must be interpreted together.

\medterm{Uremic Encephalopathy} Brain dysfunction caused by the buildup of waste products in severe kidney failure. It may cause confusion, extreme sleepiness, muscle twitching, seizures, or coma. It is a medical emergency and usually requires urgent treatment of kidney failure, often including dialysis.

\medterm{Uremic Frost} Rare white, powdery crystals of urea that appear on the skin when severe untreated kidney failure causes very high waste levels in the blood. It indicates advanced illness requiring urgent medical assessment and treatment of the kidney failure.

\medterm{Uremic Pericarditis} Inflammation of the sac around the heart caused by severe kidney failure and buildup of waste products. It may cause sharp chest pain, a scratchy heart sound, or fluid around the heart; urgent kidney treatment is needed because tamponade can occur.

\medterm{Uremic Platelet Dysfunction} Reduced ability of platelets to stick together and stop bleeding because of advanced kidney failure. It can cause easy bruising, nosebleeds, heavy menstrual bleeding, or bleeding from the stomach or bowel and may improve when the kidney failure is treated.

\medterm{Uremic Stomatitis} Oral mucosal inflammation occurring in advanced kidney failure, presenting as white
plaques, ulceration, or bleeding of the gums and tongue due to accumulation of uremic
toxins.

\medterm{Ureter} The paired muscular tubes (25-30 cm) conveying urine from the renal pelvis to the
urinary bladder, crossing the pelvic brim at the pelvic inlet. Three anatomical
narrowings (pelviureteric junction, pelvic brim, vesicoureteric junction) are common
sites of ureteric stone impaction.

\medterm{Ureter Duplex} A congenital duplication of the ureter, complete or partial, draining one kidney. It may
be asymptomatic or associated with vesicoureteral reflux, ureterocele, obstruction,
recurrent urinary infection, or impaired renal drainage.

\medterm{Ureter Fissus} A congenital longitudinal splitting or duplication anomaly of the ureter. Its clinical importance depends on the anatomy and whether it causes reflux, obstruction, infection, or impaired drainage.

\medterm{Ureter Obstruction} Blockage of urine flow from a kidney to the bladder, commonly caused by a stone, stricture, tumor, or external compression and potentially leading to hydronephrosis, infection, or kidney injury.

\medterm{Ureteral Calculi} Stones that form in or pass through a ureter, the tube carrying urine from a kidney to the bladder.

\medterm{Ureteral Catheters} Thin tubes passed into the ureter to drain urine, inject contrast, measure pressure, or guide procedures. Complications include infection, bleeding, perforation, migration, and obstruction.

\medterm{Ureteral Obstruction} Blockage of urine flow causing proximal dilation and increased pressure. Acute causes
renal colic with flank pain. Chronic leads to hydronephrosis and CKD. Emergent
decompression with stent or nephrostomy for infected obstructed kidney. Diagnosis by CT
or ultrasound showing hydronephrosis.

\synonyms
Blocked Ureter

\medterm{Ureteral Perforation} A tear or hole in the ureter, sometimes caused by a stone, trauma, surgery, or an endoscopic
procedure. Urine may leak into surrounding tissues, causing pain, infection, or inflammation;
prompt imaging and urologic assessment are required.

\medterm{Ureteral Polyp} Usually a benign fibroepithelial growth arising from the ureter. It may cause hematuria, flank pain, or obstruction and must be distinguished from urothelial malignancy when evaluated.

\medterm{Ureteral Retrograde Brush Biopsy} A brush passed through a ureteroscope to collect cells from the lining of the ureter or renal pelvis for cytologic evaluation of suspected urothelial cancer.

\medterm{Ureteral Stenosis} Narrowing of the ureter caused by scar tissue, prior surgery, radiation, stones, inflammation, or congenital disease, impairing urine drainage and possibly causing hydronephrosis or recurrent infection.

\medterm{Ureteral Stent} A thin tube placed inside the ureter to keep urine flowing from the kidney to the bladder when a stone, narrowing, tumor, or swelling blocks the passage. It can cause urinary urgency, discomfort, blood in the urine, infection, or mineral buildup and is usually temporary.

\medterm{Ureteral Stent (Double-J Stent)} A tube placed in the ureter from the renal pelvis to the bladder to maintain patency,
with pigtail coils at both ends (hence double-J). Indications: ureteral obstruction
(stones, stricture, malignancy), after ureteroscopy (prevent post-op obstruction), and
as a pre-treatment for complex ureteral stones.

\synonyms
Double-J Stent

\medterm{Ureteral Stricture} Narrowing of the ureter caused by fibrosis, surgery, instrumentation, radiation,
infection, malignancy, or congenital disease. It can produce upstream hydronephrosis,
pain, infection, and declining renal function.

\medterm{Uretereroenterostomy} A urinary diversion procedure connecting a ureter to a segment of intestine; this spelling is a source variant of \textit{ureteroenterostomy}.

\medterm{Ureteric Perforation} A tear or hole in the ureter that can allow urine to leak into surrounding tissue, usually after instrumentation, surgery, trauma, or a stone; it needs prompt urologic assessment.

\medterm{Ureteritis} Inflammation of the ureter, usually related to infection, obstruction, stone passage,
instrumentation, or adjacent inflammatory disease. It may cause flank or abdominal pain,
urinary symptoms, hematuria, and fever.

\medterm{Ureterocele} A balloon-like dilation of the lower ureter where it enters the bladder, usually congenital; it can obstruct urine flow and cause infection, hydronephrosis, reflux, stones, or impaired renal function.

\medterm{Ureterolithiasis} The presence of kidney stones (calculi) within the ureter, causing ureteral obstruction
and renal colic. Stones form from calcium oxalate/phosphate (most), uric acid, struvite
(infection), and cystine.

\medterm{Ureteropelvic Junction Obstruction} Impaired drainage from the renal pelvis into the proximal ureter, caused by intrinsic
narrowing, aberrant vessels, scarring, or functional peristaltic abnormality. It may
cause hydronephrosis, flank pain, infection, stones, or loss of renal function.

\medterm{Ureteroscope} A slender endoscope used to inspect and treat the ureter and, when advanced, the renal pelvis.

\medterm{Ureteroscopes} Small endoscopes used to inspect the ureter and renal collecting system and to perform procedures such as stone removal, biopsy, or stent placement.

\medterm{Ureteroscopy} Ureteroscopy uses a small endoscope passed through the urinary tract to inspect, biopsy, or treat ureteral or kidney stones and other lesions.
the urethra, bladder, and into the ureter or renal collecting system. Used primarily for
the diagnostic evaluation and therapeutic management of ureteral and renal calculi,
ureteral strictures, and upper tract urothelial carcinoma (UTUC).

\medterm{Ureterostomy Procedure} Surgery that creates an opening from a ureter to the skin or another urinary conduit to divert urine when normal drainage to the bladder is impossible or must be bypassed.

\medterm{Urethan} An older name for ethyl carbamate, a chemical compound with carcinogenic potential in experimental and exposure settings. It is not a routine therapeutic medicine.

\medterm{Urethra} The tube conveying urine from the bladder to the exterior. Male urethra (18-20 cm):
prostatic, membranous, and spongy (penile) parts; carries urine and semen; longer and
more tortuous than female. Female urethra (3-4 cm): short, straight, predisposing to
ascending urinary tract infections.

\medterm{Urethra Disorder} A condition affecting the urethra, such as infection, inflammation, stricture, trauma, congenital abnormality, or tumor, that may cause pain, discharge, weak stream, or retention.

\medterm{Urethra, Diseases Of And Injury To} Diseases or injuries affecting the urethra, including infection, narrowing, blockage, tears, or other structural problems.

\medterm{Urethracele} A herniation or prolapse of the urethra into the anterior vaginal wall, often occurring
together with a cystocele (bladder prolapse). Urethracele is caused by weakening of the
pubocervical fascia, often from childbirth, aging, or increased intra-abdominal
pressure.

\medterm{Urethral Carcinoma} A rare cancer arising from the urethra that may cause bleeding, a lump, discharge, pain, or difficult urination; diagnosis requires examination, imaging, and biopsy.

\medterm{Urethral Catheter} A tube inserted through the urethra into the bladder to drain urine, obtain a specimen, instill medication, or monitor output. Urinary infection and urethral trauma are important risks.


\medterm{Urethral Diverticulum} A periurethral outpouching that communicates with the urethral lumen, usually arising
from infected or obstructed periurethral glands. It may cause dysuria, post-void
dribbling, recurrent urinary infection, dyspareunia, or a palpable anterior-vaginal-wall
mass.

\medterm{Urethral Fistula} An abnormal channel between the urethra and another epithelial surface, often the vagina or skin. It can cause continuous urine leakage, recurrent infection, and local irritation.

\medterm{Urethral Hypermobility} Excessive movement of the urethra when abdominal pressure rises, contributing to urine leakage during coughing, lifting, or exercise.

\medterm{Urethral Infection} Infection involving the urethra, commonly caused by bacteria, including sexually transmitted organisms.
Symptoms may include burning urination, discharge, frequency, or pelvic discomfort; testing and
appropriate treatment help prevent complications and transmission.

\medterm{Urethral Meatus} The external opening of the urethra through which urine leaves the body. Its position may be abnormal in conditions such as hypospadias or epispadias.

\medterm{Urethral Meatus, Female} The external opening of the female urethra, located in the vestibule anterior to the vaginal opening. Its position and appearance may be relevant in urinary symptoms, congenital variation, or examination.

\medterm{Urethral Obstruction} Partial or complete blockage of urine leaving the bladder, often from stricture, stone, prostate disease, tumor, or a congenital valve; acute retention requires prompt treatment.

\medterm{Urethral Pressure Profile} A test that measures pressure along the urethra, the tube carrying urine out of the body. A small catheter is used to assess how well the urethra closes and may help investigate some types of urinary leakage or difficulty emptying the bladder.

\medterm{Urethral Pressures} Measurements of pressure within or along the urethra used to assess urethral closure and urinary continence.

\medterm{Urethral Sling} Synthetic or fascial tape supporting urethra for stress incontinence. Mid-urethral sling
(TVT, TOT) standard. Minimally invasive. Complications: erosion, urgency, retention.

\medterm{Urethral Sphincter} The muscular mechanism that helps maintain urinary continence by compressing the urethra. It includes smooth and striated muscle components controlled by autonomic and voluntary pathways.

\medterm{Urethral Stenosis} Scar-related narrowing of the urethral channel that restricts urinary flow and can cause spraying, weak stream, recurrent infection, retention, or bladder damage.

\medterm{Urethral Stricture} Narrowing of the urethra caused by scar tissue. It may produce a weak urinary stream, spraying, straining, incomplete
emptying, infections, or urinary retention and can result from injury, instrumentation,
infection, or inflammatory disease.

\medterm{Urethral Syndrome} Urinary frequency, urgency, dysuria, or discomfort without a clear bacterial urinary-tract infection; causes may include irritation, pelvic-floor dysfunction, or inflammation.

\medterm{Urethritis} Urethritis is inflammation of the urethra, commonly from sexually transmitted infection, causing dysuria, irritation, or discharge and requiring evaluation of partners and complications.
(Chlamydia trachomatis, Neisseria gonorrhoeae) or non-infectious causes. Presents with
dysuria, urethral discharge, and meatal itching. Diagnosis: urine nucleic acid
amplification testing (NAAT).

\medterm{Urethrocele} Bulging or downward displacement of the urethra, usually into the front wall of the vagina. It may cause pressure, urinary leakage, difficulty emptying the bladder, discomfort, or no symptoms.

\medterm{Urethrography} Contrast imaging of urethra. Evaluates urethral strictures, trauma, anomalies. Contrast imaging of the urethra, usually with X-rays. It can show narrowing, leaks, injury, fistulas, or congenital abnormalities.
\medterm{Urethroplasty} Surgery that rebuilds or widens the urine tube, most often to treat a narrowing called a stricture. It may use nearby tissue or a graft and aims to improve urine flow.

\medterm{Urethroscopy} Endoscopic examination of the urethra using a urethroscope, sometimes combined with treatment of a stricture or other lesion.

\medterm{Urge Urinary Incontinence} Involuntary leakage of urine, significantly affecting psychological, social, and physical well-being. See urinary incontinence.
\medterm{Urgency} A sudden, compelling desire to urinate that is difficult to defer—a lower urinary tract
symptom. Causes: overactive bladder (detrusor overactivity, often idiopathic), urinary
tract infection (urgency + dysuria), bladder outlet obstruction (BPH), neurogenic
bladder, interstitial cystitis, and pelvic floor disorders.

\medterm{URI} URI usually means upper respiratory infection, commonly a viral infection of the nose, throat, or upper airways causing congestion, cough, and sore throat.

\medterm{Uric Acid} A waste product formed when the body breaks down purines. High blood levels can cause gout or kidney stones, although many people with high levels have no symptoms.
\medterm{Uric Acid  - Blood} A blood test measuring uric acid, a purine-breakdown product. The result helps assess gout, kidney stones, kidney function, or effects of some cancer treatments.
\medterm{Uric Acid Nephropathy} Acute kidney injury from uric acid crystal precipitation in tubules. Occurs in tumor
lysis syndrome and severe hyperuricemia. Prevent with hydration and rasburicase.
Alkalinization of urine may be harmful as uric acid precipitates at alkaline pH unlike
other stone types.

\medterm{Uric Acid Stones} Kidney stones composed primarily of uric acid. They form more readily in acidic urine and are
associated with hyperuricemia, gout, high purine intake, or chronic diarrhea. Some can dissolve when urine is alkalinized.

\medterm{Uric Acid Urine Test} A urine test measuring uric acid excretion, used in evaluation of gout, kidney stones, uric-acid overproduction, and selected metabolic disorders.

\medterm{Uricosuric Drug} A medicine that increases renal excretion of uric acid, such as probenecid, used in selected patients with gout.

\medterm{Uridine} A nucleoside composed of uracil linked to ribose. It is a building block of RNA and a precursor of uridine nucleotides involved in carbohydrate and membrane metabolism.

\medterm{Uridine Diphosphate} A nucleotide consisting of uridine and two phosphate groups. UDP participates in glycogen metabolism and transfers activated sugars during glycosylation reactions.

\medterm{Uridine Monophosphate} A nucleotide consisting of uridine and one phosphate group. UMP is a component of RNA and an intermediate in pyrimidine-nucleotide metabolism.

\medterm{Uridine Triacetate} An oral prodrug that supplies uridine and is used as an emergency antidote for severe or life-threatening toxicity from fluorouracil or related fluoropyrimidines, under specialist direction.

\medterm{Uridine Triphosphate} A high-energy uridine nucleotide used in RNA synthesis and in the formation of activated sugar donors for glycogen and glycoprotein production.

\medterm{Urinal} A portable container used to collect urine from someone who cannot easily reach a toilet. It should be emptied, cleaned, and stored safely to reduce spills, odor, and infection risk.

\medterm{Urinalysis} Urinalysis examines urine appearance, chemistry, and sediment for evidence of infection, kidney disease, bleeding, diabetes, dehydration, or other systemic conditions.
(bacteria), protein, glucose, ketones, bilirubin, urobilinogen, blood, and hemoglobin.

\medterm{Urinary} Relating to urine or to the organs and processes involved in making, storing, and passing urine, including the kidneys, ureters, bladder, and urethra.
\medterm{Urinary Bladder} A hollow, distensible muscular organ located in the pelvic cavity posterior to the pubic
symphysis, serving as a reservoir for urine.

\medterm{Urinary Bladder Disorder} A disease affecting bladder storage or emptying, including infection, overactivity, pain syndrome, stones, tumor, neurologic dysfunction, or outlet obstruction.

\medterm{Urinary Bladder, Diseases Of} Disorders of the bladder, including infection, stones, tumors, neurogenic dysfunction, inflammation, obstruction, and impaired emptying.

\medterm{Urinary Calculi} Stones formed from crystallized minerals in the urinary tract. They may cause colicky pain, hematuria, obstruction, infection, or kidney injury depending on their size and location.

\medterm{Urinary Casts} Small tube-shaped particles formed inside kidney tubules and passed in urine. Their appearance can provide clues about kidney inflammation, bleeding, infection, or damage.

\medterm{Urinary Catheter} A hollow, flexible tube inserted through the urethra into the bladder to drain urine; it may be used intermittently or left in place.

\medterm{Urinary Catheterization} Insertion of a tube into the bladder to drain urine, measure output, obtain a specimen, or deliver treatment; infection, trauma, and blockage are possible complications.

\medterm{Urinary Catheters} Urinary catheters drain the bladder when a person cannot empty it normally or accurate urine measurement is needed, with risks including infection, trauma, blockage, and discomfort.

\medterm{Urinary Diversion} Surgical or catheter-based rerouting of urine when the normal passage from the kidneys to the urethra cannot be used, often to a stoma or bowel reservoir.

\medterm{Urinary Fistula} An abnormal connection between part of the urinary tract and another organ or the skin, allowing urine to leak outside its normal pathway. Causes include surgery, childbirth injury, inflammation, and cancer.

\medterm{Urinary Frequency} The need to urinate more often than usual, which may result from infection, bladder irritation, increased urine production, or other urinary disorders.

\medterm{Urinary Hesitation} Difficulty starting or maintaining urination, which may result from prostate enlargement, urethral narrowing, medications, nerve disorders, pain, or impaired bladder contraction.

\medterm{Urinary Incontinence} Unwanted leakage of urine. It may happen with coughing or exercise, a sudden strong urge, an overfull bladder, difficulty reaching the toilet, or nerve and muscle problems; assessment can identify treatable causes and suitable exercises, medicines, devices, or surgery.

\synonyms
Bladder Control, Loss Of
Incontinence, Urinary
Loss Of Bladder Control

\medterm{Urinary Incontinence Products} Urinary-incontinence products such as pads, protective underwear, catheters, and mattress covers manage leakage but do not treat its underlying cause.

\medterm{Urinary Lithiasis} The formation or presence of stones anywhere in the urinary tract, including the kidneys, ureters, bladder, or urethra.

\medterm{Urinary Retention} Inability to completely empty the bladder, which can be acute (sudden, painful,
inability to void) or chronic (gradual, often painless, with overflow incontinence).
Causes: benign prostatic hyperplasia, urethral stricture, neurogenic bladder,
medications (anticholinergics), and pelvic organ prolapse.

\medterm{Urinary Retraining} Scheduled toilet visits and bladder-control techniques used to reduce urgency, frequent urination, or leakage. A clinician may combine it with pelvic-floor exercises, fluid changes, or treatment of the underlying cause.
\medterm{Urinary Stasis} Slowing or stopping of normal urine flow, which can promote urinary infection, stone formation, or obstruction-related kidney injury.

\medterm{Urinary System} The kidneys, ureters, bladder, and urethra, which produce, transport, store, and eliminate urine while regulating fluid, electrolytes, acid--base balance, and blood pressure.

\medterm{Urinary Tract} The kidneys, ureters, bladder, and urethra. Together these organs make urine, carry it to the bladder, store it, and remove it from the body.
\medterm{Urinary Tract Disorder} Any disease of the kidneys, ureters, bladder, or urethra, including infection, stones, obstruction, congenital abnormalities, cancer, and functional or neurologic disorders.

\medterm{Urinary Tract Infection} Infection of the urinary tract, classified as lower infection such as cystitis or urethritis,
or upper infection such as pyelonephritis. Symptoms vary; dysuria, frequency, urgency, fever,
flank pain, or nausea may occur. Escherichia coli is a common cause, but organisms vary by
setting and patient factors.

\synonyms
Bladder Infection (Urinary Tract Infection)

\medterm{Urinary Tract Infection In Pregnancy} A urinary infection or asymptomatic bacterial growth occurring during
pregnancy. Pregnancy increases the risk of upper urinary infection, so screening and treatment follow pregnancy-specific
clinical guidance.

\medterm{Urinary Tract Neoplasm} A benign or malignant growth in the kidneys, ureters, bladder, or urethra; evaluation depends on its site, symptoms, tissue type, and extent of spread.

\medterm{Urinary Urgency} A sudden, difficult-to-delay need to urinate. It may occur with urinary infection, bladder irritation,
overactive bladder, stones, neurologic disease, or other conditions. Urgency with fever, flank
pain, blood in the urine, or inability to urinate requires medical assessment.

\medterm{Urination} The biological process of releasing urine from the bladder through the urethra out of
the body; also known as micturition.

\synonyms
micturition.

\medterm{Urination Disorder} Abnormal storage or emptying of urine causing frequency, urgency, dysuria, incontinence, hesitancy, weak stream, retention, or excessive or reduced urine output.

\medterm{Urination, Difficulty} Trouble starting, maintaining, or emptying urine, often described as hesitancy, weak stream, straining, or incomplete emptying; causes include obstruction, infection, medicines, and nerve or muscle disorders.

\medterm{Urine} The liquid waste produced by the kidneys, containing water and dissolved substances eliminated from the blood.

\medterm{Urine 24-Hour Volume} Twenty-four-hour urine volume is the total urine produced in a day and helps assess hydration, kidney function, endocrine disease, or abnormal urine production.

\medterm{Urine Albuminuria} An abnormal amount of the blood protein albumin in urine. Persistent albumin loss can be an early sign of kidney disease and also increases the risk of worsening kidney function and cardiovascular disease.

\medterm{Urine Alkalinization} Deliberate increase of urine pH, usually with intravenous sodium bicarbonate, to enhance elimination of selected weak acids such as salicylate; it requires monitoring of pH, potassium, fluid status, and acid-base balance.

\medterm{Urine Analysis} Laboratory examination of urine for cells, blood, proteins, glucose, microorganisms, and other substances. Laboratory testing of urine for cells, blood, protein, sugar, microbes, and other clues to illness.
\medterm{Urine Chemistry} Measurement of chemical constituents in urine, such as glucose, protein, electrolytes, ketones, blood, and pH, to evaluate kidney function and systemic disease.

\medterm{Urine Concentration Test} A test of the kidney’s ability to concentrate urine, assessed by urine osmolality or specific gravity after controlled fluid conditions when appropriate.

\medterm{Urine Conductivity} Urine conductivity measures how readily urine carries electrical current and reflects its dissolved electrolyte and solute concentration.

\medterm{Urine Culture} A laboratory test that grows microorganisms from a urine specimen to identify urinary-tract infection and guide antimicrobial selection. Results may report the organism, quantity of growth, and antimicrobial susceptibility. Proper collection is important because contamination can produce misleading growth; results must be interpreted with symptoms and urinalysis.

\medterm{Urine Culture - Catheterized Specimen} A urine culture performed on a specimen collected through urinary catheterization, usually when a clean-catch sample cannot be obtained or when catheter-associated infection is being assessed. Collection technique and timing affect interpretation because colonization and contamination may cause bacterial growth without symptomatic infection.

\medterm{Urine Cytology} Microscopic examination of cells shed into urine, mainly used to help detect high-grade urothelial carcinoma and some other abnormalities. A negative result does not exclude all urinary-tract cancers.

\medterm{Urine Drainage Bags} Urine drainage bags collect urine from a catheter or urinary diversion and require correct positioning, emptying, hygiene, and monitoring for blockage or infection.

\medterm{Urine Formation} The process by which the kidneys make urine: they filter blood, return needed water and substances to the body, and add selected wastes and extra ions to the fluid that leaves as urine.

\medterm{Urine HCG (Pregnancy Test)} Rapid point-of-care test detecting human chorionic gonadotropin in urine. Sensitivity:
detects as low as 25 mIU/mL. Positive as early as 12-14 days after conception (before
missed period). Results in 3-5 minutes. False negative: dilute urine, very early
pregnancy, ectopic pregnancy.

\synonyms
Pregnancy Test

\medterm{Urine Hemoglobin} Hemoglobin detected in urine, which may reflect intravascular hemolysis or contamination with intact red blood cells. Interpretation requires comparison with urine microscopy and clinical findings.

\medterm{Urine Osmolality} Concentration of urine relative to plasma normally 50 to 1200 mOsm/kg. High greater than
500 in prerenal states reflecting intact concentrating ability. Low less than 350
indicates impaired concentrating ability in diabetes insipidus, chronic kidney disease,
and excessive fluid intake. More precise than specific gravity.

\medterm{Urine PH} A measurement of urine acidity or alkalinity. It varies with diet, metabolism, medications, infection, and kidney function and can influence the formation of some urinary stones.

\medterm{Urine PH Test} A test measuring urine acidity or alkalinity, useful in evaluating acid-base handling, urinary infection, kidney stones, and some metabolic conditions.

\medterm{Urine Protein Dipstick Test} A screening test using a reagent strip to detect mainly albumin in urine; positive results require confirmation and quantitative assessment when clinically indicated.

\medterm{Urine Protein Electrophoresis Test} A laboratory test that separates urine proteins by electrical charge and size to help detect abnormal proteins, including monoclonal light chains associated with plasma-cell disorders.

\medterm{Urine Retention} Inability to empty the urinary bladder completely or at all. Acute retention can cause severe pain and requires prompt medical assessment.

\medterm{Urine Sediment} Microscopic examination of centrifuged urine providing diagnostic information about
renal pathology. Normal findings include few RBCs, WBCs, and hyaline casts.

\medterm{Urine Sodium} Concentration of sodium in urine measured in mEq/L. Low less than 20 in prerenal states
reflecting avid tubular reabsorption. High greater than 40 in intrinsic renal disease,
diuretic use, and salt-wasting nephropathies. Should be interpreted in context of volume
status, medications, and concurrent therapies.

\medterm{Urine Specific Gravity Test} A test estimating urine density relative to water and reflecting solute concentration; results help assess hydration and renal concentrating ability.

\medterm{Urine Toxicology} Laboratory analysis of urine for medicines, drugs, or their metabolites. Results require interpretation because detection does not
always establish impairment or the timing of exposure.

\medterm{Urinometer} A hydrometer used to measure the specific gravity of urine.

\medterm{Urobilinogen} A colourless product formed in the intestine from bilirubin and partly reabsorbed and excreted in urine; abnormal levels may provide clues to liver disease or haemolysis.


\medterm{Urochordata} A subphylum of chordates, also called tunicates, whose members include sea squirts. They are primarily biological research organisms rather than human pathogens.

\medterm{Urodynamic Testing} Suite of diagnostic procedures designed to evaluate the physiological function of the
lower urinary tract (bladder and urethral sphincter) during filling, storage, and
emptying phases.

\medterm{Urodynamics} A set of tests that assess the function of the lower urinary tract (bladder and
urethra).

\medterm{Uroflow} A test that measures how quickly and steadily urine leaves the body during urination. The result can suggest blockage, weak bladder muscles, or other emptying problems, but it is interpreted with symptoms and sometimes additional tests.

\medterm{Uroflowmetry} A noninvasive test measuring the rate and pattern of urine flow during voiding to evaluate obstruction, weak detrusor function, or voiding dysfunction.
and testing conditions; a low or abnormal flow may reflect obstruction, weak bladder
contraction, or inadequate effort and requires clinical correlation.

\medterm{Urogenital} Relating to the urinary and reproductive systems or structures involving both. The term may describe organs, infections, cancers, birth conditions, injuries, symptoms, tests, or treatments affecting these closely connected body systems.

\medterm{Urogenital Fistula} An abnormal connection between the urinary and genital tracts, such as a vesicovaginal or ureterovaginal fistula. It commonly causes continuous urinary leakage and requires specialist evaluation.

\medterm{Urogenital System} The related urinary and reproductive organs and their associated ducts, vessels, and supporting tissues. The systems share developmental origins and anatomic relationships.

\medterm{Urogenital Tract} The combined urinary and reproductive organs and passages. These include kidneys, urine tubes, bladder, urethra, reproductive organs, and connected structures that develop partly from related tissues before birth.

\medterm{Urography} Imaging of the urinary tract, traditionally using X-rays after contrast dye is given. It can show the kidneys, urine tubes, bladder, narrowing, stones, leaks, blockage, or other structural problems.

\medterm{Urokinase} A plasminogen activator produced by the body and also used as a thrombolytic medicine to help dissolve blood clots.

\medterm{Urolithiasis} Formation of stones in the urinary tract, including the kidneys, ureters, bladder, or urethra. Stone composition
varies by population and may include calcium oxalate, calcium phosphate, uric acid,
struvite, or cystine.

\medterm{Urologic Diagnostic Yield} The proportion of urologic diagnostic procedures or tests that provide a definitive and
clinically useful result.

\medterm{Urologist} A doctor who specializes in the urinary system and, in many practices, the male reproductive system. Urologists diagnose and treat kidney stones, infections, urine leakage, prostate problems, cancers, infertility, and conditions requiring procedures or surgery.

\medterm{Urology} Surgical and medical specialty concerned with the study, diagnosis, and management of
diseases affecting the male and female urinary tract (kidneys, ureters, urinary bladder,
urethra) and the male reproductive system (prostate, testes, epididymis, seminal
vesicles, penis).

\medterm{Urosepsis} Sepsis caused by infection of the urinary tract or urinary system, often associated with pyelonephritis, obstruction, catheterization, or an infected urinary calculus.
The infection may progress from a localized urinary source to bacteremia, systemic inflammation, organ dysfunction, and septic shock.
Clinical findings may include fever or hypothermia, flank pain, dysuria, hypotension, altered mental status, oliguria, or elevated lactate.
Diagnosis combines clinical assessment with urine and blood cultures, laboratory testing, and imaging when obstruction or a complicated source is suspected.
\medterm{Urostomy Leak} Escape of urine beneath or around a urostomy appliance or from a surgical connection. Persistent leakage can irritate the skin and may signal poor appliance fit, obstruction, infection, or a surgical complication.

\medterm{Urostomy Procedure} Surgery that diverts urine through a stoma into an external pouch or conduit, commonly after bladder removal or when the lower urinary tract cannot safely transport urine.

\medterm{Urostomy Prolapse} Protrusion of a segment of bowel used as a urostomy stoma beyond its usual level. It may cause bleeding, swelling, appliance problems, or impaired urine drainage and should be assessed clinically.

\medterm{Urostomy Stenosis} Narrowing of a urinary stoma or conduit that impedes drainage, causing reduced urine output, leakage, infection, or kidney-pressure changes and sometimes requiring revision.

\medterm{Urothelial Carcinoma} Urothelial carcinoma, previously called transitional cell carcinoma, is the most common
malignancy of the urinary tract, arising from the urothelium that lines the renal
pelvis, ureters, bladder, and proximal urethra.

\medterm{Urothelium} A specialized stretchable lining of the renal pelvis, ureters, bladder, and part of the urethra. It forms a protective barrier between urine and underlying tissues and allows the bladder to expand.
\medterm{Ursodeoxycholic Acid} A hydrophilic bile acid used for primary biliary cholangitis and selected cholestatic liver diseases; it may also dissolve certain small cholesterol gallstones when bile flow is preserved.

\medterm{Urticaria} A skin reaction causing transient, itchy raised welts that change location and usually resolve within 24 hours.
It may occur with angioedema, which is deeper swelling often affecting the lips, eyelids,
tongue, or extremities.

\medterm{Urticaria Multiforme} A hypersensitivity reaction in children, often triggered by infections or medications,
characterised by widespread, annular or polycyclic erythematous plaques with dusky,
ecchymotic, or targetoid centres and pseudovesication, differing from true urticaria by
the presence of dusky discolouration and prolonged lesion duration.

\medterm{Urticaria Pigmentosa} A cutaneous form of mastocytosis in which excessive mast cells accumulate in the skin,
producing tan or brown macules or papules that may swell and itch when rubbed. Flushing,
abdominal symptoms, or systemic involvement can occur.

\medterm{Urticarial Vasculitis} Inflammation of small blood vessels associated with hive-like lesions that often last longer than ordinary urticaria and may burn, bruise, or leave residual discoloration.

\medterm{Usher Syndrome} An inherited disorder combining sensorineural hearing loss with retinitis pigmentosa,
the most common cause of combined deaf-blindness. Three types are distinguished by
severity and age of onset; management requires coordinated audiological,
ophthalmological, and genetic care.

\medterm{Ustekinumab} Ustekinumab is a monoclonal antibody blocking interleukin-12 and interleukin-23 signaling, used for psoriasis, psoriatic arthritis, inflammatory bowel disease, and other approved conditions; infection risk and hypersensitivity require monitoring.

\medterm{Ustilago} A genus of smut fungi that infect plants, especially grasses and cereals. Human infection is rare and generally occurs as an opportunistic infection in people with impaired immunity.

\medterm{Usual Interstitial Pneumonia} A pattern of permanent scarring in the lungs, often seen in idiopathic pulmonary fibrosis. Scarring is usually greatest near the outer and lower parts of the lungs and can form small spaces called honeycombing. It causes increasing breathlessness and cough and is assessed with lung tests and CT imaging.

\medterm{Ut Dict.} An abbreviation of the Latin phrase “ut dictum,” meaning “as stated” or “as directed.” In medical notes it refers the reader to information already given or to an instruction previously specified.

\medterm{Uterine (Fallopian) Tubes} Paired muscular tubes extending from the uterus toward the ovaries. They transport the oocyte and are the usual site of fertilization; blockage can contribute to infertility or ectopic pregnancy.

\synonyms
fallopian tubes; uterine tubes

\medterm{Uterine Adnexae} The uterine adnexa are structures beside the uterus, chiefly the fallopian tubes, ovaries, and supporting tissues.

\medterm{Uterine Atony} Failure of the uterus to contract adequately after delivery, resulting in postpartum
hemorrhage. The most common cause of PPH. Risk factors include prolonged labor, multiple
gestation, macrosomia, oxytocin augmentation, and uterine fibroids. Diagnosis is
clinical: boggy, soft uterus with excessive bleeding.

\medterm{Uterine Balloon Tamponade} A treatment for severe bleeding after childbirth when medicines do not control bleeding. A soft balloon is placed inside the uterus and filled with sterile fluid to press against the uterine wall and slow or stop blood loss while other treatment is arranged.

\medterm{Uterine Bleeding} Uterine bleeding is bleeding from the uterus and may result from menstruation, pregnancy complications, fibroids, polyps, hormonal disorders, cancer, or medicines.

\medterm{Uterine Cancer} Cancer arising in the uterus, most commonly from the endometrium. Abnormal vaginal bleeding,
bleeding after menopause, pelvic pain, or unusual discharge requires evaluation; early diagnosis
often permits effective treatment.

\medterm{Uterine Cavity} The uterine cavity is the internal space of the uterus lined by endometrium and is where implantation and pregnancy occur.

\medterm{Uterine Contraction} Tightening of the uterine muscle, occurring normally during menstruation and labor and sometimes abnormally with pain, pregnancy complications, or medication effects.

\medterm{Uterine Fibroid (Leiomyoma)} A benign tumor arising from the smooth muscle of the uterus. Fibroids may be
submucosal, intramural, or subserosal and may cause heavy menstrual bleeding, pelvic
pressure, pain, urinary symptoms, or reproductive problems.

\synonyms
Leiomyoma

\medterm{Uterine Fibroid Embolization} A minimally invasive procedure in which particles block blood flow to uterine fibroids, causing them to shrink. It can reduce heavy bleeding and pressure symptoms but may cause pain, infection, or effects on fertility.
\medterm{Uterine Fibroids} Benign tumors of uterine smooth muscle, also called leiomyomas. They may cause heavy menstrual bleeding, pelvic pressure, pain, infertility, or no symptoms.

\synonyms
leiomyomas.
Polyps Uterus (Uterine Fibroids)
Tumors Uterine (Uterine Fibroids)
Womb Growths (Uterine Fibroids)

\medterm{Uterine Fistula} A uterine fistula is an abnormal passage connecting the uterus with another organ or the skin, often after surgery, childbirth injury, infection, or cancer.

\medterm{Uterine Fornix} The uterine or vaginal fornix is a recess around the cervix where the upper vagina surrounds the projecting lower uterus.

\medterm{Uterine Gland} A small gland in the lining of the uterus that produces fluid and helps support the early embryo. Its activity changes during the menstrual cycle and pregnancy.

\medterm{Uterine Growths} Uterine growths may include fibroids, polyps, adenomyosis, or rarely cancer. They can cause heavy or irregular bleeding, pelvic pressure, pain, infertility, or no symptoms; examination and imaging determine the type and treatment.

\medterm{Uterine Hemorrhage} Heavy or abnormal bleeding from the uterus, occurring during pregnancy, after childbirth, or at other
times. Causes include pregnancy complications, fibroids, hormonal disorders, and malignancy;
severe bleeding with dizziness or weakness is an emergency.

\medterm{Uterine Inversion} A rare, life-threatening complication of childbirth in which the uterus turns partly or completely inside out and may protrude into or beyond the vagina. It can cause severe bleeding, shock, and intense pain and requires immediate obstetric treatment to restore the uterus and control blood loss.

\medterm{Uterine Leiomyosarcoma} A rare malignant tumor of uterine smooth muscle that may cause abnormal bleeding, pelvic pain, or an enlarging mass; diagnosis is based on tissue examination and staging.

\medterm{Uterine Monitoring} Assessment of uterine contractions and, in pregnancy, their frequency, duration, and strength to evaluate labor progress, induction response, or excessive uterine activity.

\medterm{Uterine Myomectomy} Surgical removal of uterine fibroids while preserving the uterus; it may be performed through hysteroscopy, laparoscopy, or open surgery depending on fibroid size and location.

\medterm{Uterine Obstruction} Blockage that prevents normal passage through the uterus or cervix, which may cause menstrual outflow obstruction, infertility, pregnancy complications, or retained blood or tissue.

\medterm{Uterine Perforation} A full-thickness breach of the uterine wall, most commonly occurring during
instrumentation such as dilation and curettage or hysteroscopy, potentially causing
hemorrhage and requiring surgical repair.

\medterm{Uterine Polyps} Growths arising from the inner lining of the uterus that may cause irregular bleeding, heavy periods, or fertility problems.

\synonyms
Polyps, Endometrial
Polyps, Uterine

\medterm{Uterine Prolapse} Descent of the uterus toward or through the vaginal opening because the supporting pelvic
floor and ligaments have weakened. Risk factors include childbirth, aging, obesity, and
conditions that increase chronic pressure in the abdomen.

\synonyms
Pelvic Support Problems, Uterine Prolapse
Prolapsed Uterus

\medterm{Uterine Rupture} A full-thickness tear through the uterine wall, potentially life-threatening to both
mother and fetus. Risk factors include prior caesarean section (especially classical
incision), oxytocin augmentation, and operative vaginal delivery. Presents with sudden
abdominal pain, loss of fetal heart rate, and maternal hemorrhage.

\medterm{Uterine Sarcoma} Uterine sarcoma is a rare malignant tumor arising from uterine muscle, connective tissue, or related cells and may cause bleeding, pain, or an enlarging mass.

\medterm{Uterine Tachysystole} More than five uterine contractions in 10 minutes averaged over 30 minutes, whether or
not accompanied by fetal heart-rate abnormalities. It can reduce uteroplacental
perfusion, particularly when caused by oxytocin or prostaglandins.

\medterm{Uterine Tube} A paired muscular tube that transports the oocyte from the ovary toward the uterus and is the usual site of fertilization.

\medterm{Utero-Ovarian Ligament} A band of tissue connecting each ovary to the uterus. It helps hold the ovary in position within the pelvis and lies near the fallopian tube and other pelvic supporting tissues.

\synonyms
Ovarian Ligament; Proper Ovarian Ligament

\medterm{Uterosacral Ligament} One of two strong bands of tissue supporting the uterus and top of the vagina by attaching them toward the sacrum at the back of the pelvis. Weakness can contribute to pelvic-organ prolapse; endometriosis can also affect or cause pain in this area.

\medterm{Uterus} Hollow muscular organ in the pelvis between the bladder and rectum, approximately 7.5 cm
long in adult females, responsible for menstruation, implantation, and fetal development
during pregnancy.
\commonname{Womb}

\medterm{Uterus Didelphys} A congenital condition in which two separate uteruses develop, often with two cervixes and a dividing wall in the vagina. Many people can become pregnant, but miscarriage, premature birth, abnormal fetal position, or other pregnancy complications are more common.

\synonyms
Double Uterus; Didelphic Uterus

\medterm{Uterus Disorder} A condition affecting the uterus, including fibroids, adenomyosis, endometriosis, infection, abnormal bleeding, congenital anomalies, prolapse, and cancer.

\medterm{Uterus Neoplasm} A benign or malignant growth arising in the uterus, including fibroids, endometrial cancer, and sarcomas; diagnosis depends on examination, imaging, and tissue sampling.

\medterm{Uterus, Diseases Of} Conditions affecting the womb, including fibroids, adenomyosis, endometriosis, infection, unusual bleeding, prolapse, and cancer. Symptoms may include pain, heavy periods, infertility, pressure, discharge, or bleeding after menopause.

\medterm{UTI} Infection of the urinary tract, classified as lower (cystitis, urethritis) or upper
(pyelonephritis). See Urinary Tract Infection.

\medterm{Utricle} One of the otolith organs of the inner ear, detecting horizontal linear acceleration and head position relative to gravity.

\medterm{UV Light} Ultraviolet radiation, an invisible part of sunlight or artificial light that can affect the skin and eyes and may cause sunburn, aging,
or skin cancer.

\medterm{Unsaturated Fat} A dietary fat whose fatty-acid chains contain one or more carbon-carbon double bonds. Monounsaturated and polyunsaturated fats are common in vegetable oils, nuts, seeds, and fish; replacing saturated fat with unsaturated fat generally supports healthier LDL cholesterol levels.

\medterm{Uvea} The uvea is the vascular middle layer of the eye, comprising the iris, ciliary body, and choroid; it regulates light entry, supports the lens, and supplies the outer retina.
focus, and blood supply to the eye.

\medterm{Uveal Disease} A disorder of the iris, ciliary body, or choroid, including uveitis, tumors, vascular disease, and inflammation that may cause eye pain, photophobia, floaters, or visual loss.

\medterm{Uveal Melanoma} The most common primary intraocular malignancy in adults, arising from melanocytes
within the uveal tract (choroid, ciliary body, or iris), presenting with visual
symptoms, a pigmented intraocular mass, or incidental finding on examination.

\medterm{Uveal Neoplasm} A tumor arising in the uveal tract of the eye—the iris, ciliary body, or choroid; malignant melanoma may threaten vision and spread to distant organs.

\medterm{Uveitis} Inflammation of the uveal tract—the iris, ciliary body, and choroid—and sometimes nearby eye tissues. It may result from infection, autoimmune or other inflammatory disease, trauma, or an unknown cause. Symptoms can include eye pain, redness, light sensitivity, floaters, blurred vision, or reduced vision; prompt eye assessment is important because untreated inflammation can damage vision.

\synonyms
Chorioretinitis (Uveitis)

\medterm{Uveoparotid Fever} A manifestation of sarcoidosis characterized by the triad of bilateral parotid gland
enlargement, anterior uveitis, and fever, sometimes associated with facial nerve palsy.
It reflects granulomatous inflammation and is treated with systemic corticosteroids and
potentially disease-modifying agents.

\medterm{Uvula} A small, fleshy flap of tissue that hangs from the back of the throat over the root of
the tongue.

\medterm{Uvulitis} Uvulitis is inflammation and swelling of the uvula from infection, trauma, allergy, dehydration, or irritation and may interfere with swallowing or airway protection.
dehydration, or irritation. It may cause throat pain, difficulty swallowing, gagging, or
a foreign-body sensation; marked swelling with breathing difficulty requires urgent
assessment.

\medterm{Uvulopalatopharyngoplasty} A surgical procedure that removes or reshapes selected tissue of the soft palate, uvula, and pharynx to enlarge the upper airway. It may be used in selected patients with obstructive sleep apnea, but effectiveness and risks depend on airway anatomy and other treatments.

\medterm{Ultra-processed food} Ultra-processed food is an industrially formulated product made mostly from refined ingredients and additives, often with little intact whole food. Frequent high intake is associated with poorer health outcomes, although effects vary by the food and overall diet.

\medterm{Ultra-Processed Foods} Ultra-processed foods are packaged products made through multiple industrial processes and often containing additives, added sugars, salt, or refined fats. A balanced diet should emphasize minimally processed foods such as vegetables, fruits, legumes, and whole grains.

\medterm{Ureteric Calculi} Ureteric calculi are kidney stones that have moved into a ureter, the tube carrying urine from a kidney to the bladder. They can cause severe colicky pain, blood in the urine, nausea, or urinary obstruction.

\medterm{Uropathy} Uropathy is disease or damage affecting the urinary tract or kidneys. It may result from obstruction, infection, stones, reflux, medicines, or other conditions and can impair kidney function.

\medterm{Uterine Artery Embolization} Uterine artery embolization is a minimally invasive procedure that blocks blood flow to uterine fibroids, causing them to shrink. It may reduce bleeding and pressure symptoms, but fertility and pregnancy effects require individualized discussion.
